% ===================================================================
%  TABELL Nr. 1 — D'Schoul. — De Lycée. — D'Klass.
%
%  Ceci n'est PAS une transcription : c'est une traduction, faite en
%  2026, en luxembourgeois standard. D'ou trois differences avec
%  texto/io et texto/fr, et trois seulement :
%
%   * pas de \nl ni de \cc ;
%   * pas de \begin{VUpage} ;
%   * le gras se pose sur le TERME seul, ponctuation dehors.
%
%  Les cles « %%K » sont celles de texto/io, macro pour macro pour
%  l'apparat, et les renvois \textsuperscript{(n)} visent les MEMES
%  objets de la planche.
%
%  QUINZIEME COLONNE DES DIX-SEPT, et la plus petite langue du
%  volume : un demi-million de locuteurs, langue officielle depuis
%  1984 seulement, et qui partage son pays avec le francais et
%  l'allemand.
%
%  D'OU UNE DIFFICULTE QU'AUCUNE DES QUATORZE COLONNES PRECEDENTES
%  N'AVAIT POSEE. Partout jusqu'ici, un mot qui avait l'air etranger
%  etait suspect, et le balayage cherchait des signes venus du
%  dehors. ICI LA MOITIE DU VOCABULAIRE SCOLAIRE EST FRANCAIS, ET
%  C'EST DU LUXEMBOURGEOIS : « Cahier » le cahier, « Cours » la
%  lecon, « Lycée », « Grammaire », « Portier », « Kalkül », «
%  Modeller ». Les ecrire autrement serait ecrire faux.
%
%  ET LE DANGER EST DONC L'AUTRE VOISIN, L'ALLEMAND, QUI NE SE VOIT
%  PAS. Le luxembourgeois est un parler francique mosellan : il
%  ressemble a l'allemand mot pour mot, et une phrase allemande y
%  passerait sans qu'un signe la denonce. Le balayage de cette
%  colonne ne peut donc pas travailler sur le code du caractere,
%  comme les quatorze precedents ; il travaille sur des MOTS ENTIERS
%  — « das », « ist », « und », « nicht », « auch », « eine » — qui
%  sont allemands et que le luxembourgeois ecrit autrement : « dat »,
%  « ass », « an », « net », « och », « eng ».
%
%  Il ne reste au balayage qu'un seul controle par caractere, et il
%  est sur : L'ESZETT. Le luxembourgeois n'ecrit jamais « ß » —
%  « grouss », « Faass », « ësst » prennent deux s.
%
%  LE CABLAGE N'AJOUTE RIEN A NUMERO_TAB, et c'est le second cas du
%  genre apres le galicien : « TABELL Nr. 1 » est mot pour mot la
%  ligne suedoise, et le membre suedois la couvre. CENO, lui, recoit
%  « Zeen », qui ne ressemble ni au « Szene » allemand dont il vient
%  ni au « scena » roman.
% ===================================================================

%%K t01-apar-0 apar
\VUsaut{2.0mm}
\VUcentre{12.6pt}{120}{ERKLÄRENDE BICHELCHEN}
\VUsaut{3.2mm}
\VUcentre{8.4pt}{160}{ZU}
\VUsaut{2.6mm}
\VUtitre{15.6pt}{0}{dem Delmas senge Hëllefstabellen}
\VUsaut{3.4mm}
\VUornamento{9pt}
\VUsaut{5.0mm}
\VUcentre{13.2pt}{90}{ÉISCHT SERIE}
\VUsaut{3.0mm}
\VUfilet{16mm}
\VUsaut{5.6mm}

%%K t01-apar-1 apar
\VUtitre{15.0pt}{60}{TABELL Nr. 1}
\VUsaut{1.4mm}
\VUtitre{11.4pt}{0}{D'Schoul. --- De Lycée. --- D'Klass.}
\VUsaut{2.2mm}
\VUfilet{20mm}
\VUsaut{6.4mm}

%%K t01-tit-1 sub
\VUpk{10.2pt}{40}{Allgemeng Beschreiwung.}
\VUsaut{3.0mm}

%%K t01-01-1 p
1. --- Dës Tabell weist eng \VUgras{Klass} an engem \VUgras{Lycée}. An dëser Klass
sinn aacht \VUgras{Dëscher} \textsuperscript{(18)} an aacht \VUgras{Bänk}
\textsuperscript{(32)}. Op all Bänk sëtzen dräi Schüler. De \VUgras{Léierer}
\textsuperscript{(7)} sëtzt op engem \VUgras{Stull} \textsuperscript{(8)} op senger
\VUgras{Kathedra} \textsuperscript{(6)}.

%%K t01-02-1 p
2. --- Lénks vun der Kathedra steet e \VUgras{Glasschaf} \textsuperscript{(15)}.
Riets ass déi \VUgras{schwaarz Tafel} \textsuperscript{(1)} mat hirem
\VUgras{Läppchen} \textsuperscript{(2)} an hirer \VUgras{Kräid} \textsuperscript{(3)}.

%%K t01-02-2 p
Iwwer der Kathedra hänken \VUgras{Wandtabellen} \textsuperscript{(9, 11, 12)} an eng
\VUgras{Kaart} \textsuperscript{(10)}.

%%K t01-03-1 p
3. --- Net all Schüler sëtzen op hirer \VUgras{Plaz} \textsuperscript{(17)}. De Jang
\textsuperscript{(4)} steet un der Tafel; hien hält \VUgras{dat} Läppchen a wëscht déi
\VUgras{geometresch Figuren} ewech.

%%K t01-03-2 p
De Péiter ass bei der Kathedra; hie sammelt d'\VUgras{schrëftlech Aufgaben}
\textsuperscript{(5)} a bréngt se dem Léierer.

%%K t01-04-1 p
4. --- \VUgras{Dëse} Léierer ass de Léierer fir \VUgras{lieweg Sproochen.} Hie léiert
d'Schüler friem Sproochen schwätzen, liesen a schreiwen \VUgras{(däitsch, englesch,
spuenesch, italienesch, russesch} oder \VUgras{franséisch)}.

%%K t01-05-1 p
5. --- D'Klass ass ganz grouss a ganz hell. Hannen ass e breede \VUgras{Fënster} mat
zwee \VUgras{Fënsterlueden} \textsuperscript{(78)} an eng \VUgras{Glasdier,} fir se ze
lëften an ze beliichten.

%%K t01-05-2 p
An dëser Klass ass et net kal. An der Mëtt steet, fir se ze hëtzen, en immens gudden
\VUgras{Uewen} \textsuperscript{(46)} mat sengem \VUgras{Rouer} \textsuperscript{(45)}.

%%K t01-05-3 p
Un der Mauer sinn \VUgras{Kleederhaken} \textsuperscript{(58)} befestegt; drun hänken
d'Kappen an d'Mäntel vun de Schoulkanner.

%%K t01-tit-2 sub
\VUsaut{5.2mm}
\VUpk{10.2pt}{40}{De Schoulhaff.}
\VUsaut{3.6mm}

%%K t01-06-1 p
6. --- D'\VUgras{Auer} \textsuperscript{(63)} weist néng Auer a fënnef Minutten.

%%K t01-06-2 p
D'\VUgras{Paus} \textsuperscript{(76)} ass grad eriwwer an d'Cours fänkt nees un.
D'Paus ass am \VUgras{Haff} \textsuperscript{(75)}. Mir gesinn e puer Schüler, déi
spillen. En \VUgras{Iwwerwaacher} \textsuperscript{(74)} passt op si op.

%%K t01-07-1 p
7. --- Am Haff gesi mer nach verschidde Leit, nämlech den \VUgras{Inspekter}
\textsuperscript{(73)}, de Schoulleeder (\VUgras{den Direkter})
\textsuperscript{(71)} an den \VUgras{Zensor} \textsuperscript{(72)}.

%%K t01-07-2 p
Den Inspekter kontrolléiert d'Schoulen, den Direkter leet de Lycée, den Zensor
iwwerwaacht d'Studien.

%%K t01-07-3 p
Déi véiert Persoun ass de \VUgras{Portier} \textsuperscript{(77)}. Hien dréit ënnert
dem Aarm e Register, fir déi Feelend anzeschreiwen.

%%K t01-08-1 p
8. --- Hannen am Haff sinn d'\VUgras{Turngeräter} \textsuperscript{(70)} an d'Atelier
fir d'\VUgras{Handaarbechten} \textsuperscript{(69)}.

%%K t01-08-2 p
Riets vun der Tabell gesi mer d'\VUgras{Säll} fir \VUgras{Musek} a
\VUgras{Zeechnen}.

%%K t01-noto-1 noto
\VUnotes{143.2mm}{%
(*) Fir d'\textit{Virnimm} gëtt et och recommandéiert, si no der laténgescher Form ze
transkribéieren. Dat ass dee eenzege Wee, fir onzieleg Ofwäichungen ze vermeiden. Wéi
soll een iwwregens tëschent \textit{Giovanni, Jean, Johann, Jan, Hans, John, Ivan}
asw. wielen? Solle mir e spezielle Wierderbuch fir d'Virnimm maachen? Huele mer aus
dem Latäin hiren Nominativ a soe mer: \VUgras{Ioannes, Ioanna; Iulius, Iulia;
Iakobus; Andreas; Lukas.} (Komplett Grammaire).%
}

%%K t01-tit-3 sub
\VUpk{10.2pt}{40}{Detailléiert Beschreiwung.}
\VUsaut{2.4mm}

%%K t01-tit-4 sub
\VUpk{10.2pt}{40}{Déi gutt Schüler.}
\VUsaut{3.0mm}

%%K t01-09-1 p
9. --- D'Schüler an der éischter Bänk riets vun der Kathedra sinn den \VUgras{Henri}
(*), de \VUgras{Stiffen} an de \VUgras{Charel.}

%%K t01-09-2 p
Den Henri ass ganz opmierksam a fläisseg; hien nolauschtert dem Léierer. Op sengem
Dësch huet hien e \VUgras{Cahier} \textsuperscript{(28)}, e \VUgras{Buch}
\textsuperscript{(29)}, e \VUgras{Wierderbuch} \textsuperscript{(30)} an e
\VUgras{Federmäppchen} \textsuperscript{(31)}. Säi Buch huet vill \VUgras{Säiten}; all
Kapitel enthält e puer \VUgras{Paragraphen} an all \VUgras{Zeil} vill
\VUgras{Wierder.}

%%K t01-09-4 p
Säin Noper de Stiffen \textsuperscript{(24)} ass äifreg. Hie schreift a säi Cahier
d'Aufgab fir déi nächste Kéier. Hien huet eng schéin \VUgras{Schrëft.} Säi Buch ass
zou. Mir gesinn nëmmen den \VUgras{Deckel} \textsuperscript{(27)}. Nieft him läit säi
\VUgras{Zirkel} \textsuperscript{(26)}.

%%K t01-09-5 p
De Charel \textsuperscript{(23)} hält säi Buch an der Hand; hie mécht et op, fir seng
Aufgab nach eng Kéier ze liesen. Hien huet, wéi all Schüler, en
\VUgras{Tëntefaass} \textsuperscript{(25)} virun sech. Hien ass brav.

%%K t01-10-1 p
10. --- Um Dësch niewendrun sëtzen de \VUgras{Ludwig, den Anton} an de \VUgras{Paul.}
De Ludwig seet seng \VUgras{Lektioun} \textsuperscript{(22)} op an äntwert op
d'\VUgras{Froen} vum Léierer. Den Anton \textsuperscript{(21)} ass ganz onopmierksam;
heiansdo bestrooft de Léierer hien esouguer. De Paul \textsuperscript{(20)} ass e
gudde \VUgras{Kamerad,} frëndlech an hëllefsbereet. Hien ass ganz opmierksam.

%%K t01-tit-5 sub
\VUsaut{4.2mm}
\VUpk{10.2pt}{40}{Déi schlecht Schüler.}
\VUsaut{3.2mm}

%%K t01-11-1 p
11. --- Hannert dem Henri sëtzt de \VUgras{Jorsch} \textsuperscript{(38)}. Hien ass
kee schlechte Bouf, mä hien ass \VUgras{schwätzeg,} onopmierksam an onroueg. Seng
\VUgras{Liichtfäerdegkeet} a seng \VUgras{Ongehorsamkeet} maachen
\VUgras{Duercherneen} während der Cours, an dacks verdéngt hien eng streng
\VUgras{Ermahnung.} Säi \VUgras{Brouillon} \textsuperscript{(35)} ass dreckeg; hien
\VUgras{mécht Tëntefleken} drop mat senger \VUgras{Fieder} \textsuperscript{(34)}, an
dofir benotzt hien ëmmer säi \VUgras{Gummi} \textsuperscript{(36)}; seng
\VUgras{Schoultäsch} \textsuperscript{(33)} ass zerrass, well hien ass ganz
noléisseg. Firwat dréit hie säi \VUgras{Bläistëfthalter} \textsuperscript{(37)} an
d'Klass vun de \VUgras{liewege} Sproochen?

%%K t01-11-2 p
Säin Noper den Edmond \textsuperscript{(39)} huet hien net gär. Hien ass zwar e bësse
lues, mä hien ass roueg a rouheg. Hie schwätzt net; amplaz nozelauschteren, liest hie
léiwer d'\VUgras{Geschichten} a sengem \VUgras{Liesbuch.}

%%K t01-11-3 p
Den drëtte Schüler an dëser Bänk ass de \VUgras{Ludwig} \textsuperscript{(40)}. Hien
ass fläisseg. Hie schreift op e wäisst \VUgras{Blat Pabeier.} Virun sech huet hie säi
\VUgras{Läschpabeier} \textsuperscript{(41)} geluecht.

%%K t01-12-1 p
12. --- D'Schüler an der zweeter Bänk, lénks vum Léierer, si ganz jonk. Den
éischten, de klengen \VUgras{Emil,} kann nach net ganz gutt schreiwen; hien dréit seng
\VUgras{Schifertafel} \textsuperscript{(42)}, säi \VUgras{Schifergrëffel} a säi
\VUgras{Schwämmchen} \textsuperscript{(43)} mat.

%%K t01-12-2 p
Nieft him sinn den \VUgras{August} an de \VUgras{Bernard} \textsuperscript{(44)}. Dee
leschte schafft gutt, mä säin \VUgras{Undoen} ass ganz vernoléissegt.

%%K t01-13-1 p
13. --- Hannert hinne sëtzen dräi \VUgras{Faulenzer. Den Albert} léiert seng
Lektiounen net. De \VUgras{Jhemp} \textsuperscript{(47)} schléift amplaz
nozelauschteren. Seng \VUgras{Faulheet} bréngt him \VUgras{Virwërf} an esouguer
\VUgras{Strofen} an. Endlech ass de \VUgras{Chrëschten} \textsuperscript{(48)} ze
liichtsanneg. Hie \VUgras{behëlt} sech schlecht a bemitt sech guer net, sech ze
besseren.

%%K t01-14-1 p
14. --- Seng Noperen dogéint behuelen sech gutt. Den \VUgras{Ernest}
\textsuperscript{(51)} huet d'Uerdnung an d'Reegelméissegkeet gär; hie benotzt ëmmer
säi \VUgras{Lineal} \textsuperscript{(50)} a säi \VUgras{Bläistëft}
\textsuperscript{(49)}. Hien ass en \VUgras{externe Schüler.}

%%K t01-14-2 p
De \VUgras{Frantz} \textsuperscript{(52)} ass en \VUgras{interne Schüler}; hien dréit
eng Uniform, ësst a schléift am Lycée. Hien ass e gudde \VUgras{Kamerad.}

%%K t01-14-3 p
De Moritz spëtzt säi \VUgras{Bläistëft} mat sengem \VUgras{Messer}
\textsuperscript{(56)}; hien ass ganz suergfälteg.

%%K t01-14-4 p
Hien dréit all seng Bicher mat, seng \VUgras{Grammaire} \textsuperscript{(55)}, säi
\VUgras{Notizbuch} \textsuperscript{(54)} an esouguer seng \VUgras{Schreifënnerlag}
\textsuperscript{(53)}. Säi \VUgras{Schoulsak} \textsuperscript{(57)} läit um Buedem
hannert him.

%%K t01-14-5 p
Déi zwee Dëscher hannen an der Klass si vun de \VUgras{gudde Schüler}
\textsuperscript{(19)} besat.

%%K t01-tit-6 sub
\VUsaut{4.6mm}
\VUpk{10.2pt}{40}{Zeechnen a Musek.}
\VUsaut{3.4mm}

%%K t01-15-1 p
15. --- An der \VUgras{Zeechesall} hänken zierlech \VUgras{Modeller} un der Mauer, an
eng \VUgras{Lamp} \textsuperscript{(62)} mat engem \VUgras{Schirm} hänkt vun der
Plafong erof. De \VUgras{Léierer} \textsuperscript{(60)} ass ganz gedëlleg; hie
verbessert d'\VUgras{Zeechnungen} \textsuperscript{(61)} vun de Schüler a léiert si
eng \VUgras{Bust} \textsuperscript{(59)} no der Natur zeechnen.

%%K t01-16-1 p
16. --- De \VUgras{Museksall} schéngt ganz interessant. Do léiert een d'\VUgras{Musek}
liesen, d'\VUgras{Tounleedernoten} \textsuperscript{(67)} solfegéieren, déi op den
\VUgras{Notelinnen} geschriwwe sinn (do, re, mi, fa, sol, la,
si\,=\,C. D. E. F. G. A. B.), d'\VUgras{Schlësselen} kennen a charmant
\VUgras{Melodien} sangen. D'Museksstécker gi op de \VUgras{Pult}
\textsuperscript{(65)} geluecht. Ee vun de Schüler \VUgras{spillt Piano}
\textsuperscript{(64)} an de \VUgras{Museksléierer} \textsuperscript{(66)}
\VUgras{gëtt den Takt} \textsuperscript{(68)}.

%%K t01-17-1 p
17. --- Mir gesinn en Eck vum Atelier fir d'\VUgras{Handaarbechten}
\textsuperscript{(69)}, wou d'Jongen dierfen léieren, d'Holz ze hiwwelen an d'Eisen ze
feilen. Dat gëtt et net an all Lycée.

%%K t01-tit-7 sub
\VUsaut{5.0mm}
\VUpk{10.2pt}{40}{De Sall fir d'Naturwëssenschaften.}
\VUsaut{3.6mm}

%%K t01-18-1 p
18. --- De Mathematiksléierer léiert d'\VUgras{Schüler} \VUgras{Geometrie, Rechnen,
Kalkül} an d'\VUgras{metrescht System}. Op der Tafel gesi mer déi wichtegst
geometresch Figuren: de \VUgras{Krees} \textsuperscript{(a)}, de \VUgras{Quadrat}
\textsuperscript{(b)}, den \VUgras{Dräieck} \textsuperscript{(c)}, d'\VUgras{Linn}
\textsuperscript{(d)}.

%%K t01-18-2 p
Mir gesinn och eng \VUgras{Rechenoperatioun;} et ass keng \VUgras{Additioun,} keng
\VUgras{Subtraktioun,} keng \VUgras{Multiplikatioun}: et ass jo eng
\VUgras{Divisioun} \textsuperscript{(e)}.

%%K t01-19-1 p
19. --- Op der Tabell vum metresche System gesi mer: de \VUgras{Meter}
\textsuperscript{(a)}, d'\VUgras{Wo} \textsuperscript{(b)} a säi \VUgras{Gewiicht,} de
\VUgras{Liter} \textsuperscript{(e)}, den \VUgras{Dekaliter} \textsuperscript{(f)},
den \VUgras{Deziliter} an de \VUgras{Kubikmeter} \textsuperscript{(d)}.

%%K t01-19-2 p
D'Schüler studéieren och \VUgras{Naturwëssenschaft} \textsuperscript{(11)}, Zoologie
\textsuperscript{(a)}, Botanik \textsuperscript{(b)}, Geologie \textsuperscript{(c)}
an \VUgras{Geschicht} \textsuperscript{(12)}.

%%K t01-19-3 p
Am Schaf sinn d'Apparater fir \VUgras{Physik} \textsuperscript{(15)} a fir
\VUgras{Chimie} \textsuperscript{(16)}.

%%K t01-19-4 p
Um Schaf sinn: eng \VUgras{Kugel} \textsuperscript{(14)}, e \VUgras{Kegel} an eng
\VUgras{Äerdkugel} \textsuperscript{(13)}.

%%K t01-19-5 p
Iwwer den Tabellen hänkt eng \VUgras{Kaart,} déi déi fënnef Deeler vun der Welt
duerstellt: \VUgras{Europa} \textsuperscript{(g)}, \VUgras{Asien}
\textsuperscript{(e)}, \VUgras{Afrika} \textsuperscript{(f)}, \VUgras{Amerika}
\textsuperscript{(ab)} an \VUgras{Ozeanien} \textsuperscript{(h)}, an déi zwee grouss
Ozeanen, de \VUgras{Pazifik} \textsuperscript{(c)} an den \VUgras{Atlantik}
\textsuperscript{(d)}.
\VUsaut{6.0mm}
\VUfilet{22mm}
